\chapter{Muscle Weakness}

\section*{Definition}
Muscle weakness is a true reduction in the force a muscle can generate, and it is different from fatigue (tiredness after effort), pain-limited movement, or general sleepiness. The pattern -- which muscles are weak, how quickly the weakness started, and whether it is worse at the end of the day -- points toward which part of the nervous system or muscle is at fault.

\section*{When to See a Doctor}
Get emergency care for any of the following:
\begin{itemize}
  \item Sudden weakness on one side of the body, facial droop, or trouble speaking -- possible stroke
  \item Weakness that begins in the feet and climbs upward with numbness or tingling -- possible Guillain-Barre syndrome
  \item New difficulty breathing or swallowing
  \item Double vision or drooping eyelids with weakness
  \item Rapidly progressive weakness, weakness after a fall or injury, or loss of bladder or bowel control
\end{itemize}

\section*{How It Develops}
A muscle contraction begins as an electrical command from the motor area of your brain, travels down the spinal cord, out a peripheral nerve, and across the neuromuscular junction (the tiny gap between nerve and muscle) to trigger the muscle fiber itself. Weakness occurs when any stop along this path is disrupted: the brain (stroke), the spinal cord (compression), the peripheral nerve (neuropathy), the junction (myasthenia gravis), or the muscle (myopathy). The distribution is a clue. Proximal weakness -- trouble rising from a chair, climbing stairs, or lifting arms overhead -- suggests muscle disease or a junction problem. Distal weakness -- dropping things or a foot that drags -- suggests nerve disease. Weakness that comes and goes or worsens with repeated use points to myasthenia gravis. One-sided weakness that starts suddenly points to the brain or spinal cord.

\section*{Causes}
\begin{itemize}
  \item Brain and spinal cord: Stroke, spinal cord compression, multiple sclerosis, and cervical or lumbar nerve compression
  \item Peripheral nerve: Diabetic and other neuropathies, vitamin B12 deficiency, Guillain-Barre syndrome, and nerve injury
  \item Neuromuscular junction: Myasthenia gravis, in which weakness fluctuates and fatigues with use
  \item Muscle: Inflammatory myositis, thyroid disease, muscular dystrophy, and statin or steroid-related myopathy
  \item Metabolic: Low potassium, calcium, magnesium, or sodium
  \item Motor neuron: Amyotrophic lateral sclerosis (ALS), with combined weakness, wasting, and twitching
  \item Other: Chronic illness, bed rest, and deconditioning; alcohol and certain toxins
\end{itemize}

\section*{Related Symptoms}
\begin{itemize}
  \item Numbness, tingling, burning, or reduced sensation
  \item Muscle wasting (shrinkage), cramps, or twitching (fasciculations)
  \item Fatigue, particularly fatigue that worsens with repeated activity
  \item Pain in the muscle or along a nerve path
  \item Double vision, drooping eyelids, or difficulty swallowing or speaking
\end{itemize}

\section*{Medical Evaluation}
\begin{itemize}
  \item A careful history establishes which muscles are weak, when it started, what makes it worse, and which medicines you take.
  \item A neurologic examination tests muscle strength (graded 0--5), reflexes, sensation, coordination, and walking.
  \item Blood tests include creatine kinase (a muscle enzyme), electrolytes, thyroid, vitamin B12, glucose, and inflammatory markers.
  \item Electromyography and nerve conduction studies (EMG/NCS) tell whether the problem lies in the muscle, the nerve, or the junction.
  \item MRI or CT imaging of the brain or spine, and sometimes a muscle biopsy or blood tests for myasthenia antibodies, complete the picture.
\end{itemize}

\section*{Treatment Options}
\begin{itemize}
  \item Physical and occupational therapy build strength, preserve function, and prevent falls -- the cornerstone of most plans.
  \item Treat the cause: correct electrolyte deficiencies, control thyroid disease and diabetes, and replace vitamin B12.
  \item Immunosuppressive drugs (corticosteroids, intravenous immunoglobulin, plasma exchange) treat inflammatory myopathy, myasthenia gravis, and Guillain-Barre syndrome.
  \item Surgery or radiation may be needed for a herniated disc or spinal cord compression.
  \item Stroke is treated with clot-dissolving therapy in the emergency setting and with rehabilitation afterwards.
  \item If weakness involves the breathing muscles, noninvasive or mechanical ventilation provides respiratory support.
\end{itemize}


\section*{Self-Care and Home Management}
\begin{itemize}
  \item If weakness is due to deconditioning, begin gentle, graded exercise -- walking or light resistance -- and increase slowly to rebuild strength without injury.
  \item Eat a balanced diet with adequate protein, and ask your doctor about checking vitamin D and B12 levels.
  \item Avoid alcohol and stop any new medication (such as a statin) only with your doctor's guidance if you suspect it causes weakness.
  \item Fall-proof your home: clear walkways, use handrails and night lights, and consider a cane or walker if your balance is affected.
  \item If weakness makes chewing or swallowing hard, choose soft foods and sit upright while eating.
\end{itemize}

\section*{References}
\begin{thebibliography}{99}
\bibitem{muscle_weakness:harrison} Jameson JL, et al. \textit{Harrison's Principles of Internal Medicine}. 21st ed. McGraw-Hill; 2022.
\bibitem{muscle_weakness:emg} Preston DC, Shapiro BE. \textit{Electromyography and Neuromuscular Disorders: Clinical-Electrophysiologic Correlations}. 4th ed. Elsevier; 2021.
\bibitem{muscle_weakness:mills} Mills KR. The basics of electromyography. \textit{J Neurol Neurosurg Psychiatry}. 2005;76(suppl 2):ii32--ii35.
\bibitem{muscle_weakness:polyneuropathy} England JD, Gronseth GS, Franklin G, et al. Practice parameter: evaluation of distal symmetric polyneuropathy. \textit{Neurology}. 2009;72:177--184.
\end{thebibliography}
